% ============================================================
%  Spacetime-Superposition Complementarity (SSC)
%  Author: Max Meeuvissen
%  Date:   May 4, 2026
%  First record preprint — version 1.0
% ============================================================
\documentclass[12pt,a4paper]{article}

\usepackage[utf8]{inputenc}
\usepackage[T1]{fontenc}
\usepackage{amsmath,amssymb,amsthm}
\usepackage[margin=2.5cm]{geometry}
\usepackage[colorlinks=true,
            linkcolor=blue!70!black,
            citecolor=blue!70!black,
            urlcolor=blue!70!black]{hyperref}
\usepackage[numbers,sort&compress]{natbib}
\usepackage{setspace}
\usepackage{fancyhdr}
\usepackage{microtype}
\usepackage{xcolor}
\usepackage{booktabs}
\usepackage{enumitem}
\usepackage{mdframed}
\usepackage{titlesec}

\onehalfspacing

% ── Theorem environments ─────────────────────────────────
\newtheorem{principle}{Principle}
\newtheorem{conjecture}{Conjecture}
\newtheorem{definition}{Definition}

% ── Styled box for open problems ─────────────────────────
\newmdenv[
  backgroundcolor=gray!8,
  linecolor=gray!50,
  roundcorner=4pt,
  innertopmargin=9pt,
  innerbottommargin=9pt,
  innerleftmargin=10pt,
  innerrightmargin=10pt
]{openbox}

% ── Page style ────────────────────────────────────────────
\pagestyle{fancy}
\fancyhf{}
\rhead{\small Meeuvissen (2026) — Preprint v1.0}
\lhead{\small Spacetime-Superposition Complementarity (SSC)}
\cfoot{\small\thepage}
\renewcommand{\headrulewidth}{0.4pt}

% ── Title ────────────────────────────────────────────────
\title{%
  \textbf{Spacetime-Superposition Complementarity (SSC):}\\[0.3em]
  \textbf{A Conceptual Pre-Geometric Framework}\\[0.3em]
  \textbf{for Cosmology, Singularities, and}\\[0.3em]
  \textbf{the Quantum Measurement Problem}\\[1.2em]
  {\large\normalfont\itshape Preprint v1.0 — Conceptual First Sketch}
}

\author{%
  \textbf{Max Meeuvissen}\\[0.4em]
  \normalsize Independent Researcher\\[0.15em]
  \normalsize \texttt{maxmee18@gmail.com}\\[0.6em]
  \normalsize\itshape
  Developed in independent dialogue with Claude (Anthropic), May 2026
}

\date{%
  \normalsize May 4, 2026\\[0.4em]
  \small
  \textit{This document constitutes a first-record preprint establishing}\\
  \textit{priority and authorship of the conceptual framework described herein.}\\
  \textit{The author had no prior knowledge of the related literature}\\
  \textit{at the time the theory was developed.}
}

% ============================================================
\begin{document}
% ============================================================

\maketitle
\thispagestyle{fancy}

\begin{abstract}
\noindent
We propose a conceptual framework — the
\emph{Spacetime-Superposition Complementarity} (SSC) principle — in
which spacetime and quantum superposition are mutually exclusive,
complementary states of a single fundamental substrate.  In the
absence of spacetime, a pre-geometric \emph{superstate} obtains: a
condition in which all possible configurations co-exist
simultaneously, without spatial or temporal distinction.  Spacetime
is not the stage on which physics occurs but the mechanism by which
the superstate collapses into concrete, distinguishable reality.  We
show that this principle provides a natural language for several open
problems: (i)~the target of cosmic expansion is the pre-geometric
superstate rather than an empty void; (ii)~gravitational singularities
mark a phase transition from spacetime back into the superstate, not a
breakdown of physics; (iii)~the black hole information paradox
dissolves because information is a spacetime-dependent concept, not a
universal fundamental; (iv)~spacetime propagates from a singular
starting point into the superstate like a wavefront; what triggered
this starting point is unanswerable within the framework because
causality is itself a product of spacetime — the theory explicitly
leaves this as a foundational open question; and (v)~a complete cyclic
cosmology follows naturally.  The framework is
conceptual and has not yet been given a mathematical formulation; all
specific mathematical open problems are enumerated.  This paper serves
as a first public record of the theory and establishes authorship by
\textbf{Max Meeuvissen}.

\bigskip\noindent
\textbf{Keywords:} pre-geometry, quantum superposition, spacetime
emergence, black hole singularities, information paradox, cyclic
cosmology, quantum measurement problem, phase transition
\end{abstract}

\newpage
\tableofcontents
\newpage

% ============================================================
\section{Introduction}
\label{sec:intro}
% ============================================================

One of the deepest unresolved tensions in fundamental physics is
whether spacetime is a fundamental entity or an emergent one.  General
Relativity (GR) treats spacetime as a smooth, continuous manifold
shaped by mass-energy; Quantum Mechanics (QM) operates within a fixed
spacetime background and describes physical systems in terms of
superposition and probabilistic evolution.  The two frameworks are
mathematically inconsistent at extreme scales, and any theory of
quantum gravity must eventually resolve this tension.

The present paper introduces a conceptual framework developed
independently by the author in May 2026.  The starting point was
deliberately naive: \emph{What is ``empty'' space, fundamentally?}
Rather than invoking the quantum vacuum or a spacetime background,
the author pursued the question through step-by-step logical reasoning
and arrived at the following answer: the absence of spacetime is not
nothing.  It is a \emph{pre-geometric superstate} — a condition in
which all possible configurations exist simultaneously, unconstrained
by location, time, or causal order.

From this single conceptual shift, a cascade of implications follows
for the measurement problem, cosmic expansion, black hole singularities,
the information paradox, and the origin of the universe.

\subsection{Scope and Transparency}

This paper is explicitly a \emph{conceptual first sketch}.  It makes
no claim to mathematical completeness.  What it does claim:
\begin{enumerate}[label=(\alph*)]
  \item a clear statement of the logical structure of the theory
        (Sections~\ref{sec:superstate}–\ref{sec:cyclic});
  \item its relation to — and distinction from — the existing literature
        (Section~\ref{sec:related});
  \item an explicit enumeration of the mathematical work that remains
        (Section~\ref{sec:open});
  \item a documented first record with date, author, and development
        context.
\end{enumerate}

\noindent
The theory was developed in dialogue with the AI assistant Claude
(Anthropic, 2026).  The dialogue served as a structured sounding board
comparable to discussion with a knowledgeable interlocutor.  The
intellectual content, questions, and conclusions are those of the author.

\subsection{Independent Development}

The theory was developed without prior familiarity with the relevant
literature.  The author was not aware of Wheeler's pre-geometry program,
Loop Quantum Gravity, the ``It from Qubit'' research collaboration,
Subhash Kak's related work, or cyclic cosmological models at the time
of development.  Parallels to the existing literature — documented in
Section~\ref{sec:related} — were identified \emph{after} the theory
was formulated.  This is stated not as a boast but as a matter of
record: it establishes that the framework was arrived at independently,
which is relevant to questions of priority.

% ============================================================
\section{The Pre-Geometric Superstate}
\label{sec:superstate}
% ============================================================

\subsection{Rethinking ``Empty'' Space}

The conventional view treats mass-free space as a background medium:
in classical physics, an inert container; in quantum field theory, a
vacuum with zero-point fluctuations and virtual particle activity.
In both cases, space is defined \emph{within} spacetime.

We propose something categorically different.

\begin{definition}[Pre-Geometric Superstate]
\label{def:superstate}
The \emph{pre-geometric superstate} (hereafter: the superstate) is the
condition that obtains in the complete absence of spacetime.  In the
superstate:
\begin{itemize}
  \item All possible particle configurations, energy distributions,
        and field values exist simultaneously.
  \item There is no spatial location, temporal ordering, or causal
        structure.
  \item There is no distinguishability between configurations and
        therefore no information in any formal sense (Shannon, von
        Neumann, or otherwise).
  \item The superstate is not ``nothing.''  It is the maximal
        superposition of all possibilities.
\end{itemize}
\end{definition}

\noindent
This concept is categorically distinct from the QFT vacuum, which is
defined within a spacetime background and carries definite energy
density.  The superstate is \emph{prior to} spacetime, not within it.

\subsection{The Complementarity Principle}
\label{sec:complementarity}

The central thesis of this framework is stated as follows.

\begin{principle}[Spacetime-Superposition Complementarity (SSC)]
\label{pr:ssc}
Spacetime and quantum superposition are complementary, mutually
exclusive states of a single fundamental substrate.
\begin{itemize}
  \item \textbf{Where spacetime exists:} quantum superposition
        collapses into concrete states.  Entities acquire definite
        positions, momenta, and state identities.  Mass, energy, and
        distinguishability exist.
  \item \textbf{Where spacetime is absent:} total superposition
        obtains.  No decoherence mechanism is available.  No definite
        states exist.
\end{itemize}
\end{principle}

\noindent
This should be carefully distinguished from the standard quantum
mechanical picture.  In Copenhagen, many-worlds, and decoherence
interpretations alike, spacetime is a fixed background, and superposition
either persists or collapses \emph{within} it.  In SSC, spacetime is
not the arena within which superposition collapses — spacetime
\emph{is} the collapse mechanism.

\subsection{Spacetime as the Physical Collapse Mechanism}

The quantum measurement problem asks: what causes a quantum
superposition to yield a definite outcome?  Existing physical
approaches invoke environmental decoherence~\citep{zurek2003},
objective reduction tied to spacetime
curvature~\citep{penrose1996}, or spontaneous
localization~\citep{ghirardi1986}.  Observer-centric interpretations
relocate rather than resolve the question~\citep{bohr1935}.

SSC states a different answer:

\begin{principle}[Spacetime as Collapse]
\label{pr:collapse}
The existence of spacetime is the physical mechanism of quantum
collapse.  No observer, measurement apparatus, or environment is
required.  Wherever spacetime obtains, total superposition is
structurally impossible.
\end{principle}

\noindent
The question is thereby reframed: rather than asking what collapses a
superposition \emph{within} spacetime, we ask what gives rise to
spacetime such that superposition cannot persist within it.

% ============================================================
\section{Cosmological Implications}
\label{sec:cosmo}
% ============================================================

\subsection{The Target of Cosmic Expansion}
\label{sec:expansion}

The standard cosmological model describes the universe as expanding:
the metric of spacetime grows over time.  The question ``into what does
the universe expand?'' has no standard answer, because GR contains no
``outside'' — the question is typically treated as ill-posed.

SSC gives a definite answer.

\begin{conjecture}[Expansion as Superstate Propagation]
\label{conj:expansion}
The expansion of the universe is the propagation of the spacetime
boundary into the pre-geometric superstate.  There is no expansion
into empty space; there is no exterior void.  The boundary of the
observable universe is the active frontier between the spacetime phase
and the superstate.
\end{conjecture}

\noindent
Several consequences follow:
\begin{itemize}
  \item The expansion velocity of the universe and the propagation
        velocity of the spacetime boundary are \emph{the same
        quantity} — not two correlated processes but one.
  \item As the spacetime boundary advances, regions of the superstate
        are ``entered'' by spacetime: the total superposition of that
        region collapses into concrete configurations of mass, energy,
        and field states.
  \item The emergence of matter and energy at the cosmic scale is not
        an assumption; it is a consequence of Principle~\ref{pr:ssc}.
        Where spacetime arrives, concrete states must crystallize.
  \item The accelerating expansion observed cosmologically may reflect
        a property of the spacetime-superstate interface itself rather
        than a mysterious intrinsic energy density.  This hypothesis
        is speculative and left as an open problem
        (Section~\ref{sec:open}, Problem~4).
\end{itemize}

\subsection{Black Holes as Inverse Expansion: The Phase Transition Hypothesis}
\label{sec:blackholes}

GR predicts that at a gravitational singularity, spacetime curvature
diverges.  This is widely taken to signal a breakdown of the theory.
SSC offers a different reading.

\begin{conjecture}[Singularity as Phase Transition]
\label{conj:singularity}
The gravitational singularity is not a mathematical pathology or a
failure of physics.  It is a \emph{phase transition} from the
spacetime phase back into the pre-geometric superstate.
\end{conjecture}

\noindent
The analogy is instructive.  Water does not become ``infinitely hot''
at \(100\,^{\circ}\mathrm{C}\): it changes state.  A model that ignores
the liquid-to-gas transition would incorrectly predict diverging
temperature.  The model fails not because thermodynamics breaks down
but because it lacks the phase-transition description.  SSC proposes
an analogous situation for GR: the divergence of curvature at a
singularity is the symptom of a missing phase-transition term in the
equations, not an intrinsic failure of geometry.

Under Conjecture~\ref{conj:singularity}:
\begin{itemize}
  \item A black hole singularity is a region where mass density has
        driven the system across the spacetime-to-superstate phase
        boundary.
  \item Beyond the singularity there is no ``further place'' in a GR
        sense.  There is the superstate — the same substrate into
        which the universe also expands at its outer boundary.
  \item \textbf{Black holes and cosmic expansion are inverse
        manifestations of the same underlying mechanism.}  Expansion
        converts superstate into spacetime; a singularity converts
        spacetime back into superstate.  The two processes are
        opposite directions along the same spacetime-superstate
        boundary.
\end{itemize}

\subsection{The Information Paradox}
\label{sec:information}

Hawking's calculation~\citep{hawking1975} showed that black holes emit
thermal radiation and eventually evaporate, apparently destroying all
information about the infalling matter.  This contradicts the unitarity
of quantum mechanics and has generated decades of work including the
firewall argument~\citep{almheiri2013}, the island
formula~\citep{penington2020,almheiri2020}, and the original statement
of the paradox~\citep{hawking2005}.

SSC does not locate where the information goes.  It attacks the
premise.

\begin{principle}[Information as Spacetime Product]
\label{pr:information}
Information — in the sense of Shannon entropy, von Neumann entropy,
or any formal definition — presupposes distinguishability, temporal
ordering, and state identity.  These properties exist only within
spacetime.  Information is not a fundamental concept; it is a
derived concept that requires spacetime for its definition.
\end{principle}

\noindent
The argument:
\begin{enumerate}
  \item Information requires distinguishable states (\(A \neq B\)).
  \item Distinguishability requires that \(A\) and \(B\) differ in a
        spacetime-defined property (position, momentum, spin axis,
        etc.)\ or occupy distinct spacetime regions.
  \item In the superstate, there is no spacetime and therefore no
        distinguishability.
  \item Therefore, in the superstate, information is not
        \emph{destroyed} — the category of information becomes
        \emph{undefined}.
\end{enumerate}

The information paradox thus dissolves not by finding the information's
hiding place, but by recognizing that the conservation law
``information is preserved'' is a statement that applies within
spacetime and cannot, without independent argument, be extended to
regions that lie outside of spacetime.  Unitarity is a property of
time-evolution operators defined on Hilbert spaces within spacetime;
its applicability at a spacetime-to-superstate boundary is an open
question, not an axiom.

% ============================================================
\section{Cyclic Cosmology under SSC}
\label{sec:cyclic}
% ============================================================

The SSC framework implies a complete cyclic cosmology that requires no
external input, no prior cause, and no external space.

\subsection*{Phase 1 — The Superstate}

No space.  No time.  No causal ordering.  The superstate is a maximal
superposition of all possible configurations — including the
configuration in which spacetime begins to exist.  Since there is no
time axis, ``before'' and ``after'' have no meaning.

\subsection*{Phase 2 — Origin of Spacetime (The Big Bang)}

\textbf{Correction of an earlier formulation.}  A prior version of
SSC stated that the origination of spacetime requires no external cause
because it is ``one of the simultaneously realized configurations of
the superstate itself.''  This formulation is internally inconsistent:
it implicitly describes a process — simultaneous self-realization —
that would itself require either time or causal structure to be
meaningful.  The corrected position is as follows.

Spacetime does not arise spontaneously at every point of the superstate.
Rather, it propagates from a \emph{singular starting point} outward
into the surrounding superstate, like a wavefront.  This starting point
is the Big Bang.

What triggered the starting point?  Within SSC, this question is not
answerable.  Causation is a property of spacetime: cause precedes
effect only because a time axis orders events.  In the absence of
spacetime, there is no causal structure and no framework within which
an explanation could be formulated.  The theory deliberately leaves the
origin of the starting point open as a foundational limit.

\paragraph{Two approaches to the origin of spacetime.}

SSC accommodates two structurally distinct pictures of how spacetime
originates, with different degrees of theoretical commitment.

\textbf{Primary version (core thesis):}  A single spacetime episode
begins at one starting point and propagates outward as a wavefront into
the superstate.  The cause of the starting point cannot be addressed
within the theory.  This is the position adopted in the main body of
SSC.

\textbf{Alternative version (speculative outlook, not core thesis):}
Multiple spacetime ``bubbles'' originate independently at different
points within the superstate and each expands outward.  Where two
bubbles meet, they dock at an inter-dimensional boundary.  This variant
would explain why spacetime does not arise uniformly everywhere at once,
and opens the possibility of a multiverse structure embedded in the
superstate.  However, it is significantly more difficult to formalize,
raises additional questions about the nature of inter-bubble boundaries,
and is presented here solely as a speculative direction for future
investigation — not as part of the current framework.

This distinction in approach is structurally related to — but distinct
from — the Hartle-Hawking no-boundary proposal~\citep{hartle1983},
which also removes the question of ``what came before.''  In the
no-boundary proposal, imaginary time eliminates the boundary while
remaining within a spacetime-like framework.  In SSC, the pre-Big-Bang
state is not a modified spacetime but a genuinely pre-geometric
superstate.

\subsection*{Phase 3 — Expansion}

Spacetime propagates outward into the superstate.  At the advancing
boundary, the superstate collapses: mass, energy, and field
configurations crystallize.  The universe grows in extent.
Gravitational attraction accumulates as mass increases.

\subsection*{Phase 4 — Gravitational Dominance and Reversal}

At a critical point, the cumulative gravitational attraction of the
universe's mass exceeds the expansion drive.  The universe begins to
contract.  Whether this is consistent with the presently observed
accelerating expansion is an open question (Section~\ref{sec:open},
Problem~6); it may require a modified treatment of dark energy within
SSC.

\subsection*{Phase 5 — Universal Collapse}

The universe contracts.  Density increases universally.  As in a black
hole singularity, the density reaches the phase-transition threshold:
the system transitions from spacetime back into the superstate — but
now for the entire universe at once.

\subsection*{Phase 6 — Return to the Superstate}

Spacetime no longer exists.  The system is again in the pre-geometric
superstate.  All configurations are simultaneously present, including
the origination of a new spacetime.  The cycle is complete.

\bigskip
\noindent
This cyclic model is self-contained.  It requires no prior spacetime,
no external cause, and no external space.  The superstate is itself the
``substrate'' from which spacetime episodes emerge and to which they
return.

% ============================================================
\section{Relation to Existing Literature}
\label{sec:related}
% ============================================================

Parallels to the following research programs were identified by the
author \emph{after} the theory was developed.  We map each parallel
carefully and identify the points of departure.

\subsection{Pre-Geometry (Wheeler, Penrose)}

Wheeler introduced \emph{pre-geometry}: a substrate more primitive than
spacetime from which spacetime geometry emerges~\citep{wheeler1968}.
His ``It from Bit'' conjecture proposed that physical reality is
fundamentally informational~\citep{wheeler1990}.  Penrose's spin
network formalism~\citep{penrose1971} and twistor theory treat
spacetime as emergent from a combinatorial structure.  Recent
approaches to quantum gravity continue this program~\citep{oriti2009}.

\textbf{Relation to SSC.}  SSC is broadly consistent with the
pre-geometry paradigm: the superstate plays the role of Wheeler's
pre-geometric substrate.  The specific novelty of SSC is identifying
the pre-geometric substrate with \emph{total quantum superposition} and
specifying the mutual exclusivity of spacetime and superposition as the
defining dynamical relationship.  SSC also departs from Wheeler's ``It
from Bit'' by arguing that information is \emph{not} fundamental —
precisely the opposite conclusion.

\subsection{``It from Qubit'' and Holographic Approaches}

Van Raamsdonk showed that spacetime connectivity is related to quantum
entanglement between boundary degrees of
freedom~\citep{vanraamsdonk2010}.  Maldacena and Susskind's ER=EPR
conjecture identifies Einstein-Rosen bridges with EPR
entanglement~\citep{maldacena2013}.  The ``It from Qubit'' program
broadly investigates how spacetime geometry emerges from quantum
information~\citep{preskill2017}.

\textbf{Relation to SSC.}  These programs treat quantum information as
fundamental and derive spacetime from it.  SSC takes the opposite
stance: information is a spacetime product and is not defined outside
spacetime.  Additionally, SSC's superstate is not a quantum information
structure — it is a condition that \emph{precedes} the existence of
information as a defined concept.

\subsection{Loop Quantum Gravity}

Loop Quantum Gravity (LQG) quantizes spacetime directly, yielding a
discrete quantum geometry (spin foams / spin networks) at the Planck
scale~\citep{rovelli1988}.  Spacetime is not fundamental in LQG; it
emerges from the underlying combinatorial structure.  LQG
actively investigates singularity resolution, showing that the Big
Bang and black hole singularities may be replaced by quantum
bounces~\citep{ashtekar2006}.

\textbf{Relation to SSC.}  LQG and SSC share the view that spacetime
is not fundamental.  LQG replaces continuous spacetime with discrete
quantum geometry.  SSC replaces it with a pre-geometric superposition
substrate.  SSC's phase-transition hypothesis for singularities is at
least formally compatible with LQG's singularity-resolution picture,
though the underlying mechanisms are different.

\subsection{Subhash Kak (2018)}

Kak has written on the intersection of quantum mechanics and
cosmology~\citep{kak2018}.  His work contains the structural claim
that the loss of superposition is related to the appearance of spatial
expansion — a parallel that is notable given SSC's independent
identification of spacetime propagation with superposition collapse.

\textbf{Relation to SSC.}  This parallel was identified after the
theory was formulated.  SSC arrived at a structurally similar idea
independently.  Kak's framework differs in motivation, scope, and
formal apparatus.  A precise comparison requires a detailed reading of
the relevant papers, which is left for future work.  The author notes
that the exact bibliographic details of Kak's 2018 contribution require
verification; the reference is included as a pointer to a related line
of thought.

\subsection{Objective Collapse Theories (Penrose, GRW)}

Penrose's objective reduction (OR) proposes that superposition is
unstable in curved spacetime and collapses when the gravitational
self-energy of the superposed states reaches the Planck
energy~\citep{penrose1996}.  The GRW spontaneous-localization
model introduces stochastic collapse without an
observer~\citep{ghirardi1986}.  Both treat collapse as a physical
(non-observer-based) process within spacetime.

\textbf{Relation to SSC.}  SSC agrees that collapse is objective and
physical.  The key distinction is categorical: in OR and GRW, collapse
happens \emph{within} spacetime due to crossing a physical threshold.
In SSC, spacetime itself \emph{is} the collapse mechanism.  The causal
direction is reversed.

\subsection{Cyclic Cosmologies}

Several cyclic cosmological models exist.  The ekpyrotic / cyclic
model~\citep{steinhardt2002} involves periodic collisions of branes.
Penrose's Conformal Cyclic Cosmology (CCC)~\citep{penrose2010} treats
the far future of one aeon as the Big Bang of the next, using conformal
rescaling.  Loop Quantum Cosmology provides a bounce
mechanism~\citep{ashtekar2006}.

\textbf{Relation to SSC.}  SSC's cyclic model is unique in passing
\emph{fully outside} spacetime between cycles.  In the ekpyrotic
model, CCC, and LQC bounce, some spacetime-like structure persists
throughout; the ``reset'' occurs at the boundary of a spacetime epoch.
In SSC, the inter-cycle state is the pre-geometric superstate — a
condition without spacetime, time, or information.  The cycle is not
spacetime-to-spacetime but spacetime-to-superstate-to-spacetime.

% ============================================================
\section{What Is Genuinely New}
\label{sec:new}
% ============================================================

To the best of the author's knowledge, the following combination of
ideas does not appear in the literature in this form.  Each item is
stated as a specific claim for which the author asserts priority.

\begin{enumerate}
  \item \textbf{The SSC Principle itself}
        (Principle~\ref{pr:ssc}).  Spacetime and superposition as
        mutually exclusive, complementary states of a single
        fundamental substrate — not merely that spacetime is emergent,
        but that its existence is definitionally incompatible with
        superposition, and vice versa.

  \item \textbf{Spacetime as the physical collapse mechanism}
        (Principle~\ref{pr:collapse}).  The identification of
        spacetime's existence — not measurement, not environment, not
        a threshold condition — as the direct physical cause of
        quantum collapse.

  \item \textbf{Cosmic expansion as superstate propagation}
        (Conjecture~\ref{conj:expansion}).  The identification of
        cosmic expansion with the advance of the spacetime boundary
        into the pre-geometric superstate, such that the two are
        numerically identical — one process, not two correlated ones.

  \item \textbf{Gravitational singularity as phase transition to the
        superstate} (Conjecture~\ref{conj:singularity}).  The
        singularity as a change of state out of the spacetime phase,
        in analogy with thermodynamic phase transitions, rather than a
        breakdown of physics.

  \item \textbf{Expansion and black holes as inverse processes of the
        same mechanism.}  The outer boundary of cosmic expansion and
        the interior boundary of a black hole singularity as opposite
        directions of the spacetime-superstate interface.

  \item \textbf{Information as a spacetime product}
        (Principle~\ref{pr:information}).  The argument that
        information (Shannon, von Neumann, or otherwise) is
        constitutively dependent on spacetime and cannot be defined
        outside it — dissolving the information paradox not by
        locating the information but by challenging the universality
        of the conservation premise.

  \item \textbf{A cyclic cosmology through the superstate.}  A
        complete cosmological cycle that passes fully outside spacetime
        between episodes, requiring no prior cause and no external
        framework.
\end{enumerate}

% ============================================================
\section{Open Questions and Required Mathematical Work}
\label{sec:open}
% ============================================================

This paper is a conceptual sketch.  The following problems must be
addressed before SSC can be evaluated as a formal physical theory.
They are stated explicitly so that the boundary between the established
framework and the incomplete work is unambiguous.

\begin{openbox}
\textbf{Open Problem 1 — What triggered the starting point of spacetime?}\\
SSC identifies the Big Bang as a singular starting point from which
spacetime propagates into the superstate like a wavefront.  What caused
or triggered this starting point is explicitly unanswerable within the
framework, because causation presupposes spacetime.  Whether a more
fundamental account is possible — without invoking causal language —
remains the deepest open question of the theory.
\end{openbox}

\begin{openbox}
\textbf{Open Problem 2 — Energy accounting at the spacetime boundary.}\\
What happens energetically when the spacetime wavefront advances into
the superstate?  Does energy ``appear'' at the boundary (consistent
with conservation within spacetime) or is energy itself a spacetime-
dependent concept applicable only post-transition?  A precise formulation
requires specifying whatever structure, if any, governs the boundary.
\end{openbox}

\begin{openbox}
\textbf{Open Problem 3 — Sharpness of the transition at the expansion boundary.}\\
Is the spacetime-superstate phase transition sharp (first-order) or
continuous (second-order)?  A sharp transition implies a definite
boundary with a discontinuous change of state; a continuous transition
implies an intermediate zone with partial properties.  These scenarios
carry different observational implications for the structure of the
early universe and the cosmic horizon.
\end{openbox}

\begin{openbox}
\textbf{Open Problem 4 — Hawking radiation in SSC: does it carry information?}\\
If the gravitational singularity marks a phase transition to the
superstate, what is the SSC account of Hawking
radiation~\citep{hawking1975}?  Is it emitted from the spacetime-
superstate interface at the singularity?  SSC's
Principle~\ref{pr:information} implies it does not carry information
in the Shannon sense — but a precise derivation of Hawking radiation
within the SSC framework is required to confirm this.
\end{openbox}

\begin{openbox}
\textbf{Open Problem 5 — Dark energy within SSC.}\\
Observations indicate the cosmic expansion is accelerating.  In SSC,
expansion is the propagation of the spacetime wavefront into the
superstate.  Dark energy might reflect an intrinsic property of the
spacetime-superstate interface — a boundary tension or surface energy
term — rather than a constant energy density of spacetime itself.
A mathematical model of the interface would need to reproduce the
observed equation of state \(w \approx -1\) and to explain the
coincidence between the onset of acceleration and the current epoch.
\end{openbox}

\begin{openbox}
\textbf{Open Problem 6 — Mathematical structure of the superstate.}\\
What mathematical object describes the superstate?  Candidate
frameworks include: a universal wavefunction without a Hilbert space
defined over spacetime; a path integral over all possible geometries
(cf.\ Hartle-Hawking~\citep{hartle1983}), extended to the pre-geometric
regime; or an entirely new structure not present in the current
mathematical physics toolkit.  This is the most fundamental formal
open problem of the framework.
\end{openbox}

\begin{openbox}
\textbf{Open Problem 7 — Can the alternative (multi-bubble) approach be formalized?}\\
The speculative alternative version of SSC (Section~\ref{sec:cyclic},
Phase~2) posits that multiple spacetime bubbles originate independently
in the superstate and dock at inter-dimensional boundaries.  If this
picture is to be taken seriously, it requires: (a)~a precise definition
of ``bubble boundary'' in pre-geometric terms; (b)~a mechanism for
bubble docking that does not presuppose causal structure; and (c)~an
account of the observable consequences of a multi-bubble structure,
including potential signatures of inter-bubble boundaries in the
cosmic microwave background or large-scale structure.
\end{openbox}

% ============================================================
\section{Discussion}
\label{sec:discussion}
% ============================================================

The SSC framework was arrived at by pursuing a conceptual question
without mathematical machinery.  This is both a strength and a
limitation.

The strength lies in conceptual generality.  SSC unifies a set of
otherwise disconnected puzzles — the measurement problem, the target of
cosmic expansion, the nature of singularities, the information paradox,
the origin of the universe — under a single organizing principle.  The
internal consistency of the implications suggests the underlying
conceptual structure is productive.

The limitation is the absence of mathematical formalization.  Without
quantitative predictions, SSC cannot be directly compared to formal
theories or tested against observations.  The open problems in
Section~\ref{sec:open} are not peripheral details; they are the core
of the research programme required to transform SSC from a conceptual
framework into a physical theory.

The author notes that this trajectory — conceptual insight preceding
mathematical formalization — is not without precedent.  The equivalence
of inertial and gravitational mass, the equivalence of simultaneous
events, and Bohr's complementarity principle were all formulated
conceptually before receiving their mathematical expressions.  This is
offered not as a defense of incompleteness but as context for the
value of clearly stated conceptual frameworks.

Finally, the author underscores the importance of the following fact:
the framework was developed without knowledge of the related literature.
The post-hoc identification of parallels in pre-geometry, LQG, and
``It from Qubit'' does not diminish the originality of SSC's specific
combination of claims; rather, the convergence of independent reasoning
toward similar territory may itself be a signal that the conceptual
territory is productive.

% ============================================================
\section{Conclusion}
\label{sec:conclusion}
% ============================================================

We have presented the Spacetime-Superposition Complementarity (SSC)
framework, developed independently by Max Meeuvissen in May 2026.  The
central principle is that spacetime and quantum superposition are
mutually exclusive states of a pre-geometric substrate.  In the absence
of spacetime, a pre-geometric superstate obtains — a condition of
maximal superposition without location, time, or information.  Spacetime
is the mechanism of quantum collapse, not its arena.

From this single principle, the following results follow:
\begin{itemize}
  \item Cosmic expansion is the propagation of spacetime into the
        pre-geometric superstate.
  \item Gravitational singularities are phase transitions from
        spacetime back into the superstate.
  \item Black holes and cosmic expansion are inverse manifestations
        of the same spacetime-superstate boundary mechanism.
  \item Information is a spacetime-dependent concept; the information
        paradox rests on a false universalization of a spacetime-local
        conservation law.
  \item Spacetime propagates from a singular starting point like a
        wavefront into the superstate; the cause of this starting point
        is unanswerable within the framework, as causation presupposes
        spacetime — the theory explicitly acknowledges this limit.
  \item A complete, self-contained cyclic cosmology follows naturally.
\end{itemize}

\noindent
The framework is consistent with but distinct from pre-geometry, Loop
Quantum Gravity, and ``It from Qubit'' programs.  The specific
combination of claims enumerated in Section~\ref{sec:new} is, to the
author's knowledge, original.

This paper serves as a first public record.  Mathematical formalization
is the necessary next step and is explicitly left as future work.  The
author welcomes correspondence and collaboration toward that end.

\bigskip\bigskip

\begin{center}
\itshape
``The most incomprehensible thing about the universe\\
is that it is comprehensible.''\\[0.3em]
\upshape — Albert Einstein
\end{center}

\bigskip

\noindent\textbf{Acknowledgments.}\quad
The author thanks the AI assistant Claude (Anthropic, 2026) for
serving as a structured dialogue partner during the development of this
framework.  The intellectual content, questions, and conclusions are
those of the author.

\bigskip\noindent\textbf{Competing interests.}\quad None.

\bigskip\noindent\textbf{Funding.}\quad None.

\bigskip\noindent\textbf{First record.}\quad
This document was first written on \textbf{May 4, 2026}.  It
constitutes the author's first public statement of the theory.  The
dialogue in which the theory was developed took place on the same date.
The author asserts intellectual property of the conceptual framework
described herein under the above authorship and date.

\newpage
\bibliographystyle{unsrtnat}
\bibliography{ssc_references}

\end{document}
